\chapter{Proposal Feedback and Iteration Plan}

\section{Proposal Feedback Summary}

No panel feedback or required changes were issued at the proposal stage, so there is no panel action to answer. The committee approved the proposal on 14~September 2026 and issued no recommended modifications. The team nevertheless identified and corrected three statements of its own proposal before this report, listed in Section~3.2.

\section{Proposal Action Register}

\begin{longtable}{@{}>{\raggedright\arraybackslash}p{0.07\linewidth}>{\raggedright\arraybackslash}p{0.24\linewidth}>{\raggedright\arraybackslash}p{0.34\linewidth}>{\raggedright\arraybackslash}p{0.27\linewidth}@{}}
\caption{Proposal action register}\label{tab:panel-actions}\\
\toprule
ID & Original statement or feedback & Clarification or action & Status / evidence \\
\midrule
\endfirsthead
\toprule
ID & Original statement or feedback & Clarification or action & Status / evidence \\
\midrule
\endhead
-- & Proposal panel issued no requested changes & No mandatory correction required (14~September 2026 approval). & Not applicable; committee email. \\
T1 & Proposed embedding and hybrid retrieval & FYP-I Mid uses bounded lexical retrieval; hybrid semantic retrieval is planned for FYP-I Final. & Clarified; Sections~2.3 and~6.7; \path{docs/mid1/15_report_mid1_addendum.md}. \\
T2 & Approximately 60 airframes in proposal & Selected real-data corpus comprises 2{,}079 records from 118 aircraft. & Clarified; \path{docs/dataset_profile.md}, Section~5.7. \\
T3 & Overall nine-module product scope & Mid delivers a working vertical slice of the core modules; deferred modules remain planned. & Clarified; Chapter~1 and Chapter~9. \\
\bottomrule
\end{longtable}

\section{Iteration Strategy}

Development is iterative and incremental, as the FYP handbook prescribes \citep{fast2026}. Three practices shaped this iteration.

\begin{itemize}
\item \textbf{Ownership with a controlled interface.} Each member owns a set of files (frontend; graph, retrieval and ranking; ingest, conflict engine, CI). Shared contracts such as the diagnose request and response were frozen early, and a change to another owner's file needed an explicit hand-off decision.
\item \textbf{Pull requests with automated checks.} In this iteration work reaches main by pull request. GitHub Actions runs the backend tests (with a real Neo4j service), the frontend tests and build, and a Docker Compose smoke test on every push; the staging deployment runs only after CI succeeds.
\item \textbf{A written knowledge bank.} A \texttt{knowledge/} folder records context, flows with dependency diagrams, and short decision records for non-obvious choices. A rule file (\texttt{AGENTS.md}) requires it to be updated in the same change as the code, so design documents cannot drift from the implementation.
\end{itemize}

\section{Iteration 1 Plan}

\textbf{Objective.} Demonstrate a working maintenance diagnostic vertical slice on real historical data.

\textbf{Selected features} (from the proposal): a maintenance-focused interface; an asset history browser; a live knowledge graph; symptom-based diagnostic input; historical evidence retrieval; ranked investigation suggestions; conflict review; API integration; automated tests; GitHub Actions CI/CD.

\textbf{Planned demonstration flow:} open an asset, view its history, enter a symptom, retrieve related records, generate an evidence-backed response, display real source references, show a relevant historical correction, and show CI results.

\textbf{Deliverables and evidence.} The running application, the tests and CI runs, the evidence documents under \texttt{docs/mid1/}, and this report.

\section{Iteration 2 Plan}

FYP-1 Final targets an integrated minimum viable product: improved diagnostic ranking, hybrid (graph and semantic) retrieval, multi-turn sessions with server-side history, an asset and component graph, role-based access, a deployed demonstration, and an updated requirements and design specification with an initial algorithm evaluation. Chapter~9 gives the detailed plan, including what slipped from this iteration.

\section{Work Breakdown and Responsibilities}

\begin{table}[H]
\centering
\caption{Work packages of Iteration 1}
\label{tab:work-packages}
\small
\begin{tabularx}{\textwidth}{@{}c L >{\raggedright\arraybackslash}p{0.16\textwidth} >{\raggedright\arraybackslash}p{0.16\textwidth}@{}}
\toprule
Pkg & Work package & Owner & Status \\
\midrule
S1 & Maintenance-only navigation; isolate legacy screens & Shahmeer & Done \\
S2 & Asset list and dated history & Shahmeer & Done \\
S3 & Diagnostic workspace states and safety & Shahmeer & Done \\
S4 & Evidence navigation and review UI & Shahmeer & Done \\
S5 & Frontend tests and build quality & Shahmeer & Done; no linter yet \\
U1 & Audit of existing diagnostics & Umar & Done \\
U2 & Real-data graph retrieval (bounded) & Umar & Done \\
U3 & Deterministic fault ranking & Umar & Done \\
U4 & Citation validation and model boundary & Umar & Done \\
U5 & Diagnostic backend tests & Umar & Done \\
A1 & Data and environment baseline & Usman & Done \\
A2 & Maintenance APIs and conflict workflow & Usman & Done \\
A3 & GitHub Actions and repository health & Usman & Done \\
A4 & Backend wiring and documentation & Usman & Done \\
-- & Diagnostic Copilot, confidence/corrections, cost cascade (PRs \#5--\#7) & Usman & Done \\
-- & Knowledge graph view, CSV upload, evidence refresh (PRs \#9--\#13) & Shahmeer & Done; \#14 open at writing \\
\bottomrule
\end{tabularx}
\end{table}


Member contributions are described in Chapter~8.

\section{Progress Against Plan}

\begin{table}[H]
\centering
\caption{Planned Iteration 1 features against delivered state}
\label{tab:progress}
\small
\begin{tabularx}{\textwidth}{@{}c >{\raggedright\arraybackslash}p{0.24\textwidth} >{\raggedright\arraybackslash}p{0.16\textwidth} L@{}}
\toprule
\# & Planned feature & Status & Delivered state and evidence \\
\midrule
1 & Maintenance-focused UI & Done & Sidebar lists only maintenance screens; legacy URLs redirect \\
2 & Asset history browser & Done & 118 aircraft; dated history with FAA identifiers \\
3 & Live knowledge graph & Done & Neo4j graph; interactive 3D view of one aircraft's subgraph \\
4 & Symptom-based diagnostic input & Done & One-shot diagnose and multi-turn Copilot \\
5 & Historical evidence retrieval & Done (lexical) & Bounded retrieval; semantic retrieval not built \\
6 & Ranked investigation suggestions & Done (heuristic) & Deterministic ranking with documented weights \\
7 & Conflict review capability & Done & Review queue, corrections view \\
8 & API integration & Done & Axios client to FastAPI routes; contract tests \\
9 & Automated tests & Done & 194 backend and 171 frontend tests pass in CI \\
10 & GitHub Actions CI/CD & Done & CI on designated branches and PRs; CI-gated staging deploy \\
\bottomrule
\end{tabularx}
\end{table}


\subsection{Team pass gate}

\begin{longtable}{@{}p{0.05\linewidth}p{0.30\linewidth}p{0.14\linewidth}p{0.42\linewidth}@{}}
\caption{Mid-I pass gate and evidence}\label{tab:mid-pass-gate}\\
\toprule
ID & Gate item & Status & Verifiable evidence or qualification \\
\midrule
\endfirsthead
\toprule
ID & Gate item & Status & Evidence / qualification \\
\midrule
\endhead
A & Functional app with real asset and graph history & Met & Neo4j graph mode, live model and 2{,}079 records on staging, as checked 18~October 2026. \\
B & Diagnosis succeeds with a live model or reports a limitation & Partial & Live-model explanations worked locally; fresh staging explanations were empty on 16~October in the evaluated evidence (later attributed to a diagnose-prompt bug and fixed after \texttt{36a98a3}; see Appendix~D). API does not conceal the limitation. \\
C & Cited record opens original history & Met & Evidence deep links and relevant Assets/Diagnostics component tests. \\
D & Conflict workflow visible and testable & Met & Review queue, detected corrections and conflict-engine tests. \\
E & Tests/build execute in hosted CI & Met & GitHub Actions run \texttt{39068741787}: 194 backend tests, 171 frontend tests and real-graph Docker smoke all passed on \texttt{36a98a3}. \\
F & Evaluated revision, setup and ownership recorded & Met as documented & Git revision \texttt{36a98a3} is the evaluation anchor; setup, tests, evidence and contributions are recorded. A later packaging tag may exist; measurement remains on this SHA (Appendix~D). \\
G & Individual technical preparation & Completed (team-reported) & Members state that preparation, technical ownership and presentation readiness have been completed. Their individual implementation contributions appear in Chapter~8. \\
\bottomrule
\end{longtable}

\section{Delays, Deviations and Recovery Plan}

The commit timeline is uneven and we state it plainly. At commit \texttt{36a98a3} the repository has 115 commits (98 excluding merges), in three phases summarised in Table~\ref{tab:commits-phase}:

\begin{table}[H]
\centering
\caption{Commits by phase (counts at \texttt{36a98a3})}
\label{tab:commits-phase}
\small
\begin{tabularx}{\textwidth}{@{}>{\raggedright\arraybackslash}p{0.22\textwidth} c L@{}}
\toprule
Dates (2026) & Commits & What the work was \\
\midrule
11 Jul -- 18 Aug & 9 & Earlier software-engineering prototype of the same stack (later set aside; frozen legacy with regression tests) \\
8 -- 10 Sep & 29 & Maintenance proof of concept: domain registry, FAA acquisition, conflict engine, evaluation harness, first maintenance screens \\
15 -- 18 Oct & 77 & Mid-I integration: contract freeze, frontend rework, copilot, confidence/corrections, cost cascade, CI on main, evidence \\
\bottomrule
\end{tabularx}
\end{table}


\begin{itemize}
\item The full corpus has not been analysed. Analysis costs days of model quota on the free tier. Recovery: the cascade reduces projected requests; a bounded run on one aircraft was done and is reported in Section~7.5.1; a full run is planned only after live token measurement.
\item Authentication, training and analytics are outside this iteration by the scope decision in Table~\ref{tab:scope-modules}.
\end{itemize}
